\chapter{Future Directions (v1.3+)}
\label{chap:future-directions}

\begin{futurebox}
\textbf{Important Notice}\\
All features described in this chapter are \textbf{explicitly out of scope for v1.2}. They represent planned enhancements for v1.3 and beyond. Do not include these features in v1.2 implementation planning or resource allocation.\\
\\
This chapter is provided for:
\begin{itemize}
  \item Strategic planning and roadmap visibility
  \item Stakeholder expectation setting
  \item Architecture decisions that preserve future extensibility
\end{itemize}
\end{futurebox}

\section*{Integration with existing chapters}
\addcontentsline{toc}{section}{Integration with existing chapters}

This chapter describes future enhancements that extend the foundational Simula specification established in earlier chapters:

\begin{itemize}
  \item \textbf{Chapter~\ref{chap:extraction}} (Extraction Pipeline): Enhancement~1 
    (Multi-Modal) extends schema extraction to document images and time-series.
  \item \textbf{Chapter~\ref{chap:taxonomy}} (Taxonomy Engine): Enhancement~2 
    (Domain Calibration) adds knowledge anchors to the Algorithm~1 taxonomy builder.
  \item \textbf{Chapter~\ref{chap:generator}} (Generation Engine): Enhancement~3 
    (Closed-Loop) adds feedback integration to Algorithm~2 data synthesis.
  \item \textbf{Chapter~\ref{sec:sim-roadmap}} (Project Roadmap): The business 
    requirements in Chapter~\ref{chap:business-requirements} and implementation 
    roadmap in Chapter~\ref{chap:implementation-roadmap} are sequenced \emph{after} 
    the v1.2 milestones (M01--M07) complete.
  \item \textbf{Chapter~\ref{ch:agent-impl}} (Agent Instructions): Task S1--S6 
    remain the technical implementation tasks; the business requirements here 
    provide the governance and compliance wrapper around those tasks.
\end{itemize}

\noindent
The enhancements defined here are categorized as \textbf{Should} or \textbf{Could} 
items in the MoSCoW classification from \cref{tab:sim-moscow}, meaning they may 
be deferred to v1.3 if the Must milestones consume available capacity.

\section*{Scope of this chapter}
\addcontentsline{toc}{section}{Scope of this chapter}

This chapter defines enterprise-grade enhancements to the Simula framework, extending
the core ``mechanism design'' approach to synthetic data generation for production
deployment in regulated industries. These enhancements address the practical limitations
identified in the original research while preserving Simula's key differentiators:
dataset-level design, seedless generation, and independent control of coverage,
complexity, and quality.

\section{Enhancement Overview}

The Simula framework, as described in ``Reasoning-Driven Synthetic Data Generation
and Evaluation'' (OpenReview NALsdGEPhB, 2026) and the Google Research blog
``Designing Synthetic Datasets for the Real World,'' introduces four foundational
steps:

\begin{enumerate}
  \item \textbf{Global Diversification} --- Taxonomic sampling scaffolds for coverage
  \item \textbf{Local Diversification} --- Meta-prompt variation to prevent mode collapse
  \item \textbf{Complexification} --- Controlled difficulty progression
  \item \textbf{Quality Assurance} --- Reasoning-driven verification
\end{enumerate}

This chapter extends each step with enterprise-ready capabilities:

\begin{table}[htbp]
\centering
\caption{Simula enhancement dimensions.}
\label{tab:enhancement-dimensions}
\footnotesize
\begin{tabularx}{\textwidth}{@{}l l X@{}}
\toprule
\textbf{Enhancement} & \textbf{Simula Step} & \textbf{Enterprise Value} \\
\midrule
Multi-Modal Extension & Global + Local & Vision, time-series, tabular data support \\
Domain Adaptation Calibration & Global & Vertical-specific accuracy (finance, healthcare) \\
Closed-Loop Validation & Quality Assurance & Measurable ROI via downstream correlation \\
Programmable Workflow Engine & All Steps & Declarative specifications, audit compliance \\
\bottomrule
\end{tabularx}
\end{table}


\section{Enhancement 1: Multi-Modal Extension}

\subsection{Current Limitation}

The Simula research primarily demonstrates text-based synthetic data generation.
Enterprise use cases frequently require multi-modal data:
\begin{itemize}
  \item \textbf{Document processing}: Scanned invoices with embedded tables (Arabic AP domain)
  \item \textbf{Time-series analytics}: Financial metrics with temporal patterns (Trial Balance domain)
  \item \textbf{Structured data}: Tabular records with categorical and numerical fields
\end{itemize}

\subsection{Cross-Modal Reasoning Scaffolds}

Extend the taxonomic scaffolding approach to multi-modal scenarios by constructing
\emph{cross-modal reasoning graphs} that capture relationships between modalities.

\begin{lstlisting}[language=jsonx,caption={Cross-modal taxonomy extension.}]
@dataclass
class MultiModalTaxonomyNode(TaxonomyNode):
    """Extended taxonomy node with cross-modal relationships."""
    
    # Primary modality for this node
    primary_modality: str = "text"  # text, image, tabular, time_series
    
    # Cross-modal mappings
    cross_modal_mappings: dict[str, str] = field(default_factory=dict)
    # Example: {"image": "invoice_header_region", "text": "vendor_name_field"}
    
    # Modality-specific constraints
    modality_constraints: dict[str, Any] = field(default_factory=dict)
    # Example: {"image": {"min_resolution": [800, 600], "format": "PNG"}}

class CrossModalTaxonomyBuilder(SimulaTaxonomyBuilder):
    """Build taxonomies with cross-modal consistency."""
    
    async def build_cross_modal_taxonomy(
        self,
        factor: Factor,
        modalities: list[str],
    ) -> MultiModalTaxonomy:
        """
        Build taxonomy ensuring cross-modal consistency.
        
        For example, in medical imaging:
        disease_type -> imaging_presentation -> text_description
        ensures generated image-text pairs are semantically aligned.
        """
        # Build primary modality taxonomy
        primary_taxonomy = await self.build_taxonomy(factor)
        
        # Extend with cross-modal mappings
        for node in primary_taxonomy.get_leaf_nodes():
            mappings = await self._infer_cross_modal_mappings(
                node, modalities
            )
            node.cross_modal_mappings = mappings
        
        return MultiModalTaxonomy(
            primary_taxonomy=primary_taxonomy,
            modalities=modalities,
        )
\end{lstlisting}

\subsection{Modality-Aware Complexification}

Define modality-specific complexity metrics:

\begin{table}[htbp]
\centering
\caption{Modality-specific complexity dimensions.}
\label{tab:modality-complexity}
\footnotesize
\begin{tabularx}{\textwidth}{@{}l X@{}}
\toprule
\textbf{Modality} & \textbf{Complexity Dimensions} \\
\midrule
Text & Vocabulary diversity, syntactic complexity, reasoning depth \\
Image & Background clutter, lighting variation, occlusion, scale variance \\
Tabular & Cardinality, sparsity, categorical diversity, numerical range \\
Time-Series & Trend strength, seasonality, anomaly density, noise ratio \\
\bottomrule
\end{tabularx}
\end{table}

\subsection{SAP Integration: Document Information Extraction}

For SAP Document Information Extraction (DOX) scenarios:

\begin{lstlisting}[language=jsonx,caption={SAP DOX multi-modal integration.}]
class SAPDOXModalityAdapter:
    """Adapter for SAP Document Information Extraction."""
    
    MODALITY_MAPPING = {
        "invoice_header": ["image_region", "text_field", "table_cell"],
        "line_items": ["table_structure", "text_sequence"],
        "totals": ["numerical_field", "currency_format"],
    }
    
    def generate_cross_modal_example(
        self,
        taxonomy_path: list[str],
        complexity_level: str,
    ) -> MultiModalTrainingExample:
        """
        Generate aligned multi-modal training example.
        
        Output includes:
        - Synthetic document image (PNG/PDF)
        - Extracted text annotations
        - Structured field mappings
        - Ground truth bounding boxes
        """
        pass
\end{lstlisting}


\section{Enhancement 2: Domain Adaptation Calibration}

\subsection{Current Limitation}

Simula's seedless, reasoning-driven approach relies on the LLM's inherent knowledge
to construct taxonomies. In specialized domains (finance, healthcare, manufacturing),
this knowledge may be incomplete or biased toward common patterns, missing critical
long-tail scenarios.

\subsection{Human-in-Loop Knowledge Anchors}

Introduce lightweight ``knowledge anchors'' that guide taxonomy construction without
violating the seedless principle:

\begin{lstlisting}[language=jsonx,caption={Domain knowledge anchor injection.}]
@dataclass
class DomainKnowledgeAnchor:
    """Lightweight knowledge injection for domain calibration."""
    
    # Domain identifier
    domain: str  # "trial_balance", "arabic_ap", "regulations"
    
    # Required coverage nodes (MUST appear in taxonomy)
    required_nodes: list[str] = field(default_factory=list)
    
    # Prohibited nodes (MUST NOT appear - regulatory/compliance)
    prohibited_nodes: list[str] = field(default_factory=list)
    
    # Expert-provided complexity anchors
    complexity_anchors: dict[str, float] = field(default_factory=dict)
    # Example: {"multi_currency_reconciliation": 0.85, "simple_variance": 0.2}
    
    # Domain-specific vocabulary
    domain_vocabulary: list[str] = field(default_factory=list)

class CalibratedTaxonomyBuilder(SimulaTaxonomyBuilder):
    """Taxonomy builder with domain calibration."""
    
    async def build_calibrated_taxonomy(
        self,
        factor: Factor,
        anchors: DomainKnowledgeAnchor,
    ) -> Taxonomy:
        """
        Build taxonomy with domain expert calibration.
        
        Process:
        1. Generate initial taxonomy via LLM reasoning
        2. Validate required nodes are present
        3. Remove prohibited nodes
        4. Calibrate complexity scores against anchors
        5. Present to expert for final review (optional)
        """
        # Step 1: Initial generation
        taxonomy = await self.build_taxonomy(factor)
        
        # Step 2: Validate required nodes
        existing_nodes = {n.name.lower() for n in taxonomy.get_leaf_nodes()}
        missing = set(anchors.required_nodes) - existing_nodes
        
        if missing:
            logger.warning(f"Missing required nodes: {missing}")
            taxonomy = await self._inject_missing_nodes(taxonomy, missing)
        
        # Step 3: Remove prohibited nodes
        taxonomy = self._prune_prohibited(taxonomy, anchors.prohibited_nodes)
        
        # Step 4: Calibrate complexity
        taxonomy = self._calibrate_complexity(taxonomy, anchors.complexity_anchors)
        
        return taxonomy
\end{lstlisting}

\subsection{Vertical-Specific Complexity Metrics}

Define industry-standard complexity dimensions:

\begin{table}[htbp]
\centering
\caption{Vertical-specific complexity definitions.}
\label{tab:vertical-complexity}
\footnotesize
\begin{tabularx}{\textwidth}{@{}l l X@{}}
\toprule
\textbf{Vertical} & \textbf{Metric} & \textbf{Definition} \\
\midrule
\multirow{3}{*}{Finance} 
  & Regulatory complexity & Number of compliance rules applicable \\
  & Temporal complexity & Multi-period analysis requirements \\
  & Entity complexity & Cross-entity consolidation depth \\
\midrule
\multirow{3}{*}{Healthcare}
  & Diagnostic complexity & Differential diagnosis breadth \\
  & Protocol complexity & Multi-step treatment pathways \\
  & Comorbidity complexity & Co-occurring condition interactions \\
\midrule
\multirow{3}{*}{Manufacturing}
  & Process complexity & Operation sequence length \\
  & Tolerance complexity & Precision requirements (6-sigma) \\
  & Supply chain complexity & Multi-tier supplier dependencies \\
\bottomrule
\end{tabularx}
\end{table}

\subsection{SAP Integration: S/4HANA Domain Schemas}

\begin{lstlisting}[language=jsonx,caption={SAP S/4HANA domain calibration.}]
# Domain anchors for SAP verticals
SAP_DOMAIN_ANCHORS = {
    "trial_balance": DomainKnowledgeAnchor(
        domain="trial_balance",
        required_nodes=[
            "intercompany_elimination",
            "currency_translation",
            "statutory_adjustment",
            "management_overlay",
        ],
        prohibited_nodes=[
            "personal_account_detail",  # Privacy
            "salary_individual",        # Compliance
        ],
        complexity_anchors={
            "simple_variance_analysis": 0.2,
            "multi_currency_consolidation": 0.75,
            "complex_elimination_entry": 0.9,
        },
    ),
    "arabic_ap": DomainKnowledgeAnchor(
        domain="arabic_ap",
        required_nodes=[
            "vat_compliance",
            "withholding_tax",
            "payment_terms_discount",
            "multi_currency_invoice",
        ],
        complexity_anchors={
            "standard_invoice": 0.15,
            "credit_note_reconciliation": 0.5,
            "multi_vendor_payment_run": 0.8,
        },
    ),
}
\end{lstlisting}


\section{Enhancement 3: Closed-Loop Validation}

\subsection{Current Limitation}

Simula's quality assurance relies on reasoning-model self-evaluation. This lacks
direct feedback from downstream model performance, making it difficult to:
\begin{itemize}
  \item Quantify ROI of synthetic data generation
  \item Identify systematic blind spots in coverage
  \item Adapt generation parameters based on model learning curves
\end{itemize}

\subsection{Training-Validation Feedback Loop}

Integrate Simula with the model training pipeline to create a closed feedback loop:

\begin{lstlisting}[language=jsonx,caption={Closed-loop validation architecture.}]
@dataclass
class DownstreamPerformanceSignal:
    """Performance signal from downstream model training."""
    
    # Model identifier
    model_id: str
    
    # Performance by complexity band
    accuracy_by_complexity: dict[str, float]  # {"easy": 0.95, "hard": 0.72}
    
    # Performance by taxonomy path
    accuracy_by_taxonomy: dict[str, float]  # {"sql/aggregation/sum": 0.98}
    
    # Identified weak spots
    underperforming_nodes: list[str]
    
    # Timestamp
    evaluated_at: str

class ClosedLoopGenerator(SimulaDataGenerator):
    """Generator with downstream feedback integration."""
    
    async def generate_with_feedback(
        self,
        initial_count: int,
        downstream_evaluator: Callable,
        max_iterations: int = 3,
    ) -> list[TrainingExample]:
        """
        Generate data with iterative refinement based on downstream performance.
        
        Algorithm:
        1. Generate initial dataset
        2. Train downstream model
        3. Evaluate performance by taxonomy node
        4. Identify underperforming regions
        5. Adjust sampling weights to oversample weak regions
        6. Generate additional targeted examples
        7. Repeat until convergence or max iterations
        """
        examples = await self.generate_examples(initial_count)
        
        for iteration in range(max_iterations):
            # Train and evaluate downstream model
            signal = await downstream_evaluator(examples)
            
            # Check convergence
            if self._is_converged(signal):
                break
            
            # Adjust sampling weights
            weight_adjustments = self._compute_weight_adjustments(signal)
            self._apply_weight_adjustments(weight_adjustments)
            
            # Generate targeted examples for weak regions
            additional = await self._generate_targeted(
                signal.underperforming_nodes,
                count=initial_count // 4,
            )
            examples.extend(additional)
        
        return examples
    
    def _compute_weight_adjustments(
        self,
        signal: DownstreamPerformanceSignal,
    ) -> dict[str, float]:
        """
        Compute sampling weight adjustments based on performance gaps.
        
        Regions with lower accuracy receive higher sampling weights
        in subsequent generations.
        """
        adjustments = {}
        mean_accuracy = sum(signal.accuracy_by_taxonomy.values()) / len(signal.accuracy_by_taxonomy)
        
        for node, accuracy in signal.accuracy_by_taxonomy.items():
            if accuracy < mean_accuracy:
                # Boost underperforming regions
                gap = mean_accuracy - accuracy
                adjustments[node] = 1.0 + (gap * 2)  # Up to 2x boost
            else:
                adjustments[node] = 1.0
        
        return adjustments
\end{lstlisting}

\subsection{Adversarial Verification Layer}

Add an adversarial detector to identify systematic biases:

\begin{lstlisting}[language=jsonx,caption={Adversarial verification for distribution drift.}]
class AdversarialVerifier:
    """Detect systematic biases in synthetic data."""
    
    async def detect_distribution_drift(
        self,
        synthetic_examples: list[TrainingExample],
        reference_distribution: dict[str, float],
    ) -> DriftReport:
        """
        Compare synthetic data distribution against reference.
        
        Detects:
        - Mode collapse (over-concentration in specific patterns)
        - Coverage gaps (missing taxonomy regions)
        - Complexity skew (imbalanced difficulty distribution)
        """
        pass
    
    async def adversarial_probe(
        self,
        synthetic_examples: list[TrainingExample],
        probe_queries: list[str],
    ) -> ProbeReport:
        """
        Use adversarial probes to test for memorization.
        
        Checks if generated data inadvertently reproduces
        patterns from the reasoning model's training data.
        """
        pass
\end{lstlisting}

\subsection{SAP Integration: AI Core MLOps}

\begin{lstlisting}[language=jsonx,caption={SAP AI Core integration for closed-loop training.}]
class SAPAICoreIntegration:
    """Integration with SAP AI Core for MLOps feedback loop."""
    
    async def submit_training_job(
        self,
        training_data_path: str,
        model_config: dict,
    ) -> str:
        """Submit training job to SAP AI Core."""
        # Use SAP AI Core SDK
        pass
    
    async def get_evaluation_metrics(
        self,
        execution_id: str,
    ) -> DownstreamPerformanceSignal:
        """Retrieve evaluation metrics from completed training."""
        # Parse AI Core metrics
        pass
    
    async def trigger_retraining(
        self,
        signal: DownstreamPerformanceSignal,
        threshold: float = 0.85,
    ) -> Optional[str]:
        """Trigger retraining if performance below threshold."""
        if signal.overall_accuracy < threshold:
            # Generate additional targeted examples
            # Submit new training job
            pass
\end{lstlisting}


\section{Enhancement 4: Programmable Workflow Engine}

\subsection{Current Limitation}

Simula provides conceptual controllability but lacks:
\begin{itemize}
  \item User-friendly interfaces for business users
  \item Declarative specification of dataset requirements
  \item Version control and reproducibility guarantees
  \item Audit trails for compliance
\end{itemize}

\subsection{Declarative Dataset Specification Language}

Design a YAML-based DSL for dataset specification:

\begin{lstlisting}[language=yamlx,caption={Simula dataset specification DSL.}]
# simula-dataset-spec.yaml
apiVersion: simula/v1
kind: DatasetSpecification
metadata:
  name: trial-balance-text2sql-v3
  domain: trial_balance
  owner: finance-ai-team
  created: 2026-04-21
  
spec:
  # Target dataset size
  target_count: 10000
  
  # Coverage requirements
  coverage:
    taxonomy_paths:
      - path: "query_type/aggregation/*"
        min_coverage: 0.15
      - path: "query_type/filtering/*"
        min_coverage: 0.20
      - path: "complexity/multi_table_join"
        min_coverage: 0.10
    
  # Complexity distribution
  complexity:
    distribution:
      easy: 0.30
      medium: 0.45
      hard: 0.25
    calibration_anchors:
      - example_id: "expert_example_001"
        target_score: 0.85
  
  # Quality thresholds
  quality:
    critic_pass_rate: 0.90
    diversity_threshold: 0.70
    
  # Generation parameters
  generation:
    model: "gpt-4-turbo"
    temperature: 0.7
    best_of_n: 5
    enable_complexification: true
    complexification_ratio: 0.30
    
  # Validation
  validation:
    schema: "docs/schema/simula/training-example.schema.json"
    downstream_eval: true
    
  # Output
  output:
    format: jsonl
    path: "data/hana_generated/tb_text2sql_v3.jsonl"
    include_metadata: true
\end{lstlisting}

\subsection{Version Control and Reproducibility}

\begin{lstlisting}[language=jsonx,caption={Git-like versioning for datasets.}]
@dataclass
class DatasetVersion:
    """Version metadata for reproducibility."""
    
    version_id: str  # Semantic version or hash
    parent_version: Optional[str]
    
    # Generation configuration snapshot
    spec_snapshot: dict
    taxonomy_snapshot: dict
    
    # Reproducibility parameters
    random_seed: int
    model_version: str
    
    # Statistics
    example_count: int
    coverage_report: dict
    quality_report: dict
    
    # Audit trail
    created_by: str
    created_at: str
    approval_status: str  # draft, reviewed, approved
    approver: Optional[str]

class DatasetVersionControl:
    """Git-like version control for synthetic datasets."""
    
    def commit(
        self,
        examples: list[TrainingExample],
        spec: DatasetSpecification,
        message: str,
    ) -> DatasetVersion:
        """Commit a new dataset version."""
        pass
    
    def diff(
        self,
        version_a: str,
        version_b: str,
    ) -> DatasetDiff:
        """Compare two dataset versions."""
        pass
    
    def rollback(
        self,
        version_id: str,
    ) -> list[TrainingExample]:
        """Rollback to a previous version."""
        pass
    
    def audit_log(
        self,
        version_id: str,
    ) -> list[AuditEntry]:
        """Get complete audit trail for a version."""
        pass
\end{lstlisting}

\subsection{SAP Integration: BTP and HANA Cloud}

\begin{lstlisting}[language=jsonx,caption={SAP BTP integration for workflow orchestration.}]
class SAPBTPWorkflowIntegration:
    """Integration with SAP BTP for workflow orchestration."""
    
    def __init__(self, btp_config: BTPConfig):
        self.object_store = btp_config.object_store  # SAP Object Store
        self.hana_client = HanaClient(btp_config.hana)
        self.workflow_service = btp_config.workflow_service
    
    async def execute_workflow(
        self,
        spec: DatasetSpecification,
    ) -> WorkflowExecution:
        """
        Execute Simula workflow on SAP BTP.
        
        Steps:
        1. Parse specification
        2. Extract schema from HANA Cloud
        3. Build taxonomies
        4. Generate examples
        5. Validate against schema
        6. Upload to HANA Cloud
        7. Trigger downstream training (optional)
        """
        pass
    
    async def store_dataset_version(
        self,
        version: DatasetVersion,
        examples: list[TrainingExample],
    ) -> str:
        """Store dataset version in SAP Object Store with HANA metadata."""
        # Store examples in Object Store
        object_path = f"simula/{version.version_id}/examples.jsonl"
        
        # Store metadata in HANA Cloud
        self.hana_client.insert_dataset_version(version)
        
        return object_path
\end{lstlisting}


\section{Regulatory Framework Alignment}

These enhancements are designed to align with emerging regulatory guidance on
synthetic data:

\subsection{UK FCA Synthetic Data Expert Group}

The Financial Conduct Authority (FCA) Synthetic Data Expert Group identifies
nine governance principles\footnote{FCA, ``Synthetic Data Expert Group Report,'' 2024}:

\begin{enumerate}
  \item \textbf{Transparency} --- Enhancement 4 (audit trails, version control)
  \item \textbf{Explainability} --- Enhancement 2 (knowledge anchors, calibration)
  \item \textbf{Purpose limitation} --- Enhancement 4 (declarative specifications)
  \item \textbf{Data minimization} --- Enhancement 1 (targeted generation)
  \item \textbf{Accuracy} --- Enhancement 3 (closed-loop validation)
  \item \textbf{Security} --- SAP BTP integration (access control)
  \item \textbf{Accountability} --- Enhancement 4 (ownership, approval workflows)
  \item \textbf{Fairness} --- Enhancement 3 (adversarial bias detection)
  \item \textbf{Continuous monitoring} --- Enhancement 3 (drift detection)
\end{enumerate}

\subsection{IEEE Synthetic Data White Paper}

The IEEE white paper on synthetic data governance emphasizes\footnote{IEEE, 
``Synthetic Data: Generation, Validation, and Governance,'' 2025}:

\begin{itemize}
  \item \textbf{Privacy risk quantification} --- Enhancement 3 (adversarial probing)
  \item \textbf{Fidelity-utility tradeoff} --- Enhancement 3 (downstream validation)
  \item \textbf{Provenance tracking} --- Enhancement 4 (version control)
\end{itemize}

\subsection{World Economic Forum Guidelines}

The WEF framework for synthetic data in financial services notes\footnote{WEF, 
``Synthetic Data for Financial Services,'' 2024}:

\begin{itemize}
  \item Cross-functional governance teams are essential (addressed in Chapter~\ref{chap:business-requirements})
  \item Traceability and provenance are foundational (Enhancement 4)
  \item Continuous monitoring prevents quality degradation (Enhancement 3)
\end{itemize}


\section{Implementation Checklist}

\begin{itemize}
  \item[\texttt{[ ]}] Multi-modal taxonomy builder with cross-modal mappings
  \item[\texttt{[ ]}] Modality-specific complexification metrics
  \item[\texttt{[ ]}] SAP DOX integration for document extraction
  \item[\texttt{[ ]}] Domain knowledge anchor injection mechanism
  \item[\texttt{[ ]}] Vertical-specific complexity calibration
  \item[\texttt{[ ]}] Closed-loop training feedback integration
  \item[\texttt{[ ]}] Adversarial verification layer
  \item[\texttt{[ ]}] SAP AI Core MLOps integration
  \item[\texttt{[ ]}] Declarative DSL parser
  \item[\texttt{[ ]}] Dataset version control system
  \item[\texttt{[ ]}] SAP BTP workflow orchestration
  \item[\texttt{[ ]}] HANA Cloud metadata storage
\end{itemize}